% ECONOMICS_PROBLEM: {"id": "L42-541", "title": "Vertical integration and self-preferencing", "jel_code": "L42", "source_id": "OP-0077"}
% FIRST_POSTED: 2026-10-09
% ATTEMPT_STATUS: partial
% ENTRY_KIND: research_attempt
\documentclass[11pt,a4paper]{article}
\usepackage[T1]{fontenc}
\usepackage[utf8]{inputenc}
\usepackage{lmodern,amsmath,amssymb,amsthm,mathtools}
\usepackage[margin=25mm]{geometry}
\usepackage{microtype,enumitem,hyperref}
\hypersetup{colorlinks=true,urlcolor=blue,linkcolor=blue}
\setlength{\parindent}{0pt}
\setlength{\parskip}{4pt}
\newtheorem{theorem}{Theorem}[section]
\newtheorem{proposition}[theorem]{Proposition}
\newtheorem{lemma}[theorem]{Lemma}
\newtheorem{corollary}[theorem]{Corollary}
\newcommand{\R}{\mathbb R}
\newcommand{\E}{\mathbb E}
\newcommand{\Prob}{\mathbb P}
\newcommand{\ind}{\mathbf 1}
\DeclareMathOperator{\Var}{Var}
\DeclareMathOperator{\Cov}{Cov}
\DeclareMathOperator{\dist}{dist}
\title{Vertical efficiency and foreclosure under explicit contracts}
\author{GPT 6 Astra Ultra}
\date{9 October 2026}
\begin{document}\maketitle
\textbf{Status: Partial; the full catalogue target remains unresolved.}
\begin{abstract}
Successive linear markups create an exactly calculable integration gain, while efficient two-part contracting removes that gain before integration. In a separate competitive-input benchmark, excluding a downstream rival reduces welfare. These solved examples establish contract-dependent opposite signs. An exact path-decomposition identity explains why assigning interactions to efficiency or foreclosure requires a stated intermediate regime.
\end{abstract}
\section{A contract-sensitive target}
The catalogue asks for consumer and total welfare effects of a specified ownership and conduct regime. Ownership alone is insufficient: wholesale contracts, access obligations, ranking, entry, and investment responses determine the equilibrium. Here all demand and cost primitives are known and there is one static market. The models are intentionally restricted, with no empirical inference from observed margins.

In the first benchmark inverse final demand is $P(Q)=a-Q$ for $Q\in[0,a]$, upstream marginal cost is $c$ with $0\le c<a$, and downstream transformation is costless. Set $D=a-c>0$. Transfers between the two firms cancel in total welfare, which is
\[
 W(Q)=\int_0^Q(a-x-c)\,dx=DQ-Q^2/2.
 \tag{1}
\]
Consumer surplus at the market price is $Q^2/2$. The demand curve is a declared money-metric valuation schedule, so adding consumer surplus and joint profit reproduces (1).
\section{Successive linear markups}
There is one upstream monopolist and one downstream monopolist. The upstream firm first posts a nonnegative linear wholesale price $w$; the retailer then chooses quantity. For $w<a$, the retailer maximizes $(a-Q-w)Q$, giving $Q(w)=(a-w)/2$. The upstream firm maximizes $(w-c)(a-w)/2$, so its optimum is $w^*=(a+c)/2$. Prices above $a$ yield zero trade and cannot beat this positive-profit outcome.

After vertical integration the single firm maximizes $(a-Q-c)Q$, giving $Q_I=D/2$. The separated outcome is $Q_S=D/4$.
\begin{proposition}
In this linear-contract game, integration increases total welfare by $5D^2/32$ and consumer surplus by $3D^2/32$.
\end{proposition}
\begin{proof}
At separation, (1) gives
$W_S=D^2/4-D^2/32=7D^2/32$.
At integration, $W_I=D^2/2-D^2/8=3D^2/8$.
Subtracting yields $5D^2/32$. Consumer surplus rises from $D^2/32$ to $D^2/8$, giving $3D^2/32$. The strict concavity of both firms' one-variable objectives establishes the stated optima.
\end{proof}
With $a=10,c=2$, quantities rise from two to four. Welfare rises from fourteen to twenty-four, and consumer surplus from two to eight. These figures check the formulas with $D=8$. The efficiency gain is specific to the restriction to successive linear prices, not a general property of separate ownership.
\section{Two-part tariffs change the benchmark}
Allow an enforceable public two-part contract consisting of per-unit price $w$ and fixed fee $F$, with a retailer outside option of zero and acceptance at indifference. Set $w=c$. The retailer then chooses $Q=D/2$ and earns operating profit $D^2/4$ before the fee. Set $F=D^2/4$ to extract that profit.

This contract reproduces the integrated quantity and joint profit under separate ownership. The retailer participates and chooses the efficient-for-joint-profit quantity, while the upstream firm obtains the same total profit as an integrated monopolist. No feasible contract can generate greater joint profit than the maximum of $(a-Q-c)Q$, so the proposed contract is upstream-optimal subject to participation. Integration then has no additional output or consumer-surplus benefit in this benchmark.

A fixed fee constrained by limited liability, private demand information, hidden wholesale terms, or noncontractible investment may fail to implement this outcome. Those are substantive contract restrictions that must be measured or modeled. Observing a high linear wholesale price does not tell whether unobserved rebates or fixed payments neutralize its marginal incentive.
\section{A separate foreclosure loss}
Consider instead a competitive upstream input supplied to two identical downstream firms at marginal cost $c$. They choose quantities simultaneously under the same inverse demand. Their first-order equations are $D-2q_i-q_j=0$, giving $q_1=q_2=D/3$ and total $Q_C=2D/3$. The input market is competitive, so there is no upstream markup to eliminate.

Suppose a specified integration-and-access regime gives one downstream firm exclusive control of an essential input and bars the rival from buying it from any source. This is a genuine reduction in access, imposed as a regime primitive. The remaining firm chooses monopoly quantity $Q_F=D/2$.
\begin{proposition}
Relative to competitive input access, this exclusion lowers total welfare by $5D^2/72$ and consumer surplus by $7D^2/72$.
\end{proposition}
\begin{proof}
Equation (1) gives $W_C=2D^2/3-2D^2/9=4D^2/9$ and $W_F=3D^2/8$. Their difference is $(32-27)D^2/72$. Consumer surplus is $2D^2/9$ before exclusion and $D^2/8$ afterward, differing by $(16-9)D^2/72$.
\end{proof}
This second economy does not combine with the first by mechanically adding their welfare changes: their baseline contract and competition structures differ. Their purpose is a rigorous sign counterexample. A rival's lower profit alone would not establish exclusion; the proof specifies an essential-input access change. Efficient expansion by an integrated firm could also lower rival profits with positive consumer effects.
\section{Path dependence of mechanism accounting}
For one fully specified economy suppose the four combinations of efficiency permission $e\in\{0,1\}$ and discriminatory access $f\in\{0,1\}$ are feasible and have selected welfare $W_{ef}$. The total change is $W_{11}-W_{00}$. Assigning efficiency first gives efficiency contribution $W_{10}-W_{00}$ and access contribution $W_{11}-W_{10}$. Reversing the order gives $W_{01}-W_{00}$ and $W_{11}-W_{01}$.

The change in the attributed efficiency contribution across these orders is
\[
 I=W_{11}-W_{10}-W_{01}+W_{00}.
\]
This is immediate by subtraction. If $I\ne0$, an order-free separation into those two marginal contributions is impossible without an additional attribution rule. Averaging the two orders assigns half the interaction to each channel, but that is a convention rather than an identified causal pathway.

If one intermediate regime is legally or technologically infeasible, its welfare cannot be estimated by pretending it is an available intervention. The catalogue's $v_E$ therefore needs a real contract and access interpretation.
\section{A ranking example and the unresolved application}
For two available offers, let their net real match values be $v_O,v_R$. A ranking policy changes the probability that the consumer receives the owned offer from $p_0$ to $p_1$, holding prices, quality, participation, and costs fixed. Expected welfare changes by $(p_1-p_0)(v_O-v_R)$. The sign depends on match values, not on ownership. Allowing price or quality responses adds terms and requires another equilibrium calculation; the simple formula is only a frozen-market benchmark.

The completed proofs show why vertical welfare claims need contract data and feasible counterfactuals. They do not identify actual access changes, bargaining threats, investment, or match values in a platform. No current legal conclusion or empirical merger recommendation is made. The broad agenda remains unresolved, with the familiar markup calculations used as transparent benchmarks rather than claimed novel theory.

\section{Completion Estimate}
35\%

This is a post hoc, heuristic estimate for the full catalogue target.
The attempt computes contract-dependent integration gains and foreclosure losses, shows efficient contracting can remove successive markups, and makes mechanism path dependence explicit. Full vertical and ranking policy evaluation still requires identified access, bargaining, investment, and match-quality responses in one validated economy, with robustness across admissible conduct and equilibrium assumptions.

\begin{thebibliography}{9}
\bibitem{ht} O. Hart and J. Tirole (1990). \emph{Vertical Integration and Market Foreclosure}.
\url{https://www.brookings.edu/articles/vertical-intergration-and-market-foreclosure/}.
The primary record establishes the foreclosure-theory context; the explicit linear-demand examples here are self-contained.
\end{thebibliography}

\end{document}
